Kurs:Algebraische Kurven (Osnabrück 2025-2026)/Vorlesung 12/latex
Kurs:Algebraische Kurven (Osnabrück 2025-2026)/Vorlesung 12/latex

\setcounter{section}{12}


\epigraph { Zum Sehen geboren, Zum Schauen bestellt, } { Johann Wolfgang von Goethe }


  \zwischenueberschrift{Das $K$-Spektrum}


\bild{ \begin{center}

\includegraphics[width=5.5cm]{ \bildeinlesungjpg {Alexander_Grothendieck} {jpg} }

\end{center}

\bildtext {Alexander Grothendieck (1928-2014)} }


\bildlizenz { Alexander Grothendieck.jpg } {Konrad Jacobs} {} {Commons} {CC-BY-SA 2.0} {Quelle=Oberwolfach Photo Collection (http://owpdb.mfo.de/detail?photoID=1452)}


Wie hängen affin-algebraische Mengen und deren Koordinatenringe zusammen? Hier kann man nur für nicht-endliche Grundkörper gehaltvolle Antworten erwarten, da es im endlichen Fall zu wenige Punkte gibt. Eine befriedigende Theorie erfordert sogar, dass man sich auf algebraisch abgeschlossene Körper beschränkt, oder aber
\zusatzgs {das ist der Standpunkt der von Alexander Grothendieck entwickelten Schematheorie} {}
nicht nur $K$-Punkte betrachtet, sondern generell maximale Ideale und Primideale als Punkte mitberücksichtigt.


Eine erste wichtige Frage ist folgende: Eine $K$-Algebra $R$ von endlichem Typ hat mehrere, in aller Regel gleichberechtigte Darstellungen als Restklassenring einer Polynomalgebra, sagen wir


\mavergleichskettedisp

{\vergleichskette

{ K[X_1  , \ldots , X_n]/ {\mathfrak a}
}

{ \cong} { R
}

{ \cong} { K[X_1 , \ldots , X_m]/  {\mathfrak b}
}

{ } {
}

{ } {
}

}
{}{}{.}
Dazu gehören die beiden Nullstellengebilde


\mavergleichskette

{\vergleichskette

{ V( {\mathfrak a} )
}

{ \subseteq }{  { {\mathbb A}_{ K }^{ n  } }
}

{  }{
}

{    }{
}

{   }{
}

}
{}{}{}
und


\mavergleichskette

{\vergleichskette

{ V( {\mathfrak b} )
}

{ \subseteq }{  { {\mathbb A}_{ K }^{ m  } }
}

{  }{
}

{    }{
}

{   }{
}

}
{}{}{.}
Wie hängen diese beiden Nullstellengebilde zusammen?


\inputbeispiel{ }

{


Wir betrachten den Polynomring in einer Variablen


\mavergleichskette

{\vergleichskette

{ R
}

{ = }{ K[T]
}

{  }{
}

{    }{
}

{   }{
}

}
{}{}{.}
Ihm entspricht zunächst die affine Gerade

\mathl{{\mathbb A}^{1}_{K}}{.} Man kann $R$ aber auch auf ganz verschiedene Arten als Restklassenring einer Polynomalgebra in mehreren Variablen erhalten. Es sei beispielsweise


\mathbed {a \in K} {}

 {a \neq 0} {}

 {} {} {} {;}
wir betrachten den Restklassenring

\mathl{K[X,Y]/(aY+bX)}{.} Dieser Ring ist
\zusatzklammer {als $K$-Algebra} {} {}
isomorph zu $R$, wie die Abbildung


\mathdisp {K[X,Y]/(aY+bX) \longrightarrow K[T],\,

 \mathdisplaybruch X \longmapsto T, \, Y \longmapsto -\frac{b}{a} T} { , }


zeigt. Das zugehörige Nullstellengebilde


\mavergleichskette

{\vergleichskette

{ V(aY+bX)
}

{ \subset }{ {\mathbb A}^{2}_{K}
}

{  }{
}

{    }{
}

{   }{
}

}
{}{}{}
ist einfach die Gerade in der affinen Ebene, die durch die Gleichung


\mavergleichskette

{\vergleichskette

{ Y
}

{ = }{ - { \frac{  b }{ a } } X
}

{  }{
}

{    }{
}

{   }{
}

}
{}{}{}
beschrieben wird.


Eine weitere Möglichkeit, den Polynomring in einer Variablen als Restklassenring darzustellen, ist durch

\mathl{K[X,Y]/(Y-P(X))}{} gegeben, wobei

\mathl{P(X)}{} ein beliebiges Polynom in der einen Variablen $X$ ist. Der Ringhomomorphismus


\mathdisp {K[X,Y]/(Y-P(X)) \longrightarrow K[T],\,

 \mathdisplaybruch X \longmapsto T, \, Y \longmapsto P(T)} { , }


zeigt, dass wieder ein Isomorphismus zum Polynomring in einer Variablen vorliegt. Das zugehörige Nullstellengebilde ist einfach der
\definitionsverweis {Graph}{}{}
des Polynoms

\mathl{P(X)}{.}


\bild{ \begin{center}

\includegraphics[width=5.5cm]{ \bildeinlesungjpg {Lineline} {jpg} }

\end{center}

\bildtext {} }


\bildlizenz { Lineline.jpg } {Astur1} {} {Commons} {PD} {}


\bild{ \begin{center}

\includegraphics[width=5.5cm]{ \bildeinlesungpng {Lineair-cartesiaans} {png} }

\end{center}

\bildtext {} }


\bildlizenz { Lineair-cartesiaans.png } {MADe} {} {nl.wikipedia.org} {CC-BY-SA-3.0} {}


\bild{ \begin{center}

\includegraphics[width=5.5cm]{ \bildeinlesungpng {Polynomialdeg5} {png} }

\end{center}

\bildtext {} }


\bildlizenz { Polynomialdeg5.png } {} {Derbeth} {Commons} {CC-BY-SA-3.0} {}


}


Der Punkt an diesem Beispiel ist, dass alle drei geometrischen Objekte die Nullstellenmengen zu verschiedenen Restklassendarstellungen von

\mathl{K[T]}{} sind. Vom Standpunkt der algebraischen Geometrie sind das drei gleichberechtigte Darstellungen der affinen Geraden, auch wenn sie unterschiedlich \anfuehrung{aussehen}{.} In der algebraischen Geometrie muss man so hinschauen, dass sie gleich aussehen. Was man sieht, sind nur verschiedene Einbettungen des \anfuehrung{eigentlichen und wahren}{} geometrischen Objektes, das zu einer $K$-Algebra intrinsisch gehört, nämlich das $K$-Spektrum.


\inputdefinition

{}

{


Zu einer kommutativen
$K$-\definitionsverweis {Algebra}{}{}
$R$ von endlichem Typ bezeichnet man die Menge der
$K$-\definitionsverweis {Algebrahomomorphismen}{}{}


\mathdisp {\operatorname{Hom}_K(R,K)} {  }


als das
\definitionswortpraemath {K}{ Spektrum }{}
von $R$. Es wird mit

\mathl{K\!-\!\operatorname{Spek}\, \left(  R  \right)}{} bezeichnet.


}


Die  Elemente in einem $K$-Spektrum

\mathl{K\!-\!\operatorname{Spek}\, \left(  R  \right)}{} betrachten wir als Punkte und bezeichnen sie üblicherweise mit $P$, obwohl es definitionsgemäß Abbildungen sind, nämlich $K$-Algebrahomomorphismen von $R$ nach $K$. Für ein Ringelement


\mavergleichskette

{\vergleichskette

{ f
}

{ \in }{ R
}

{  }{
}

{    }{
}

{   }{
}

}
{}{}{}
schreiben wir dann auch einfach

\mathl{f(P)}{}
\zusatzklammer {statt \mathlk{P(f)}{}} {} {}
für den Wert von $f$ unter dem mit $P$ bezeichneten Ringhomomorphismus
\zusatzklammer {es ist nicht unüblich, einen Punkt als eine Auswertung von Funktionen anzusehen, die in einer gewissen Umgebung des Punktes definiert sind} {} {.}


Das $K$-Spektrum wird wieder mit einer \stichwort {Zariski-Topologie} {} versehen, wobei zu einem Ideal


\mavergleichskette

{\vergleichskette

{ {\mathfrak a}
}

{ \subseteq }{ R
}

{  }{
}

{    }{
}

{   }{
}

}
{}{}{}
\zusatzklammer {oder zu einer beliebigen Teilmenge aus $R$} {} {}
die Teilmenge


\mavergleichskettedisp

{\vergleichskette

{ V( {\mathfrak a} )
}

{ =} { \left\{  P \in K\!-\!\operatorname{Spek}\, \left(  R  \right)   \mid  f(P) = 0  \text{ für alle } f \in {\mathfrak a}         \right\}
}

{ } {
}

{ } {
}

{ } {
}

}
{}{}{}
als abgeschlossen erklärt wird. In der Tat wird dadurch eine Topologie definiert, siehe
Aufgabe 12.8.
Die komplementären offenen Mengen werden mit $D( {\mathfrak a} )$ bezeichnet.


\inputfaktbeweis

{Polynomring über Körper/Punkte im affinen Raum und K-Algebra-Homomorphismen/Identifizierung/Fakt}

{Lemma}

{}

{


\faktsituation {}

\faktvoraussetzung {Es sei $K$ ein
\definitionsverweis {Körper}{}{}
und sei

\mathl{K[X_1 , \ldots ,  X_n]}{} der
\definitionsverweis {Polynomring}{}{}
in $n$ Variablen.}

\faktfolgerung {Dann stehen die
$K$-\definitionsverweis {Algebrahomomorphismen}{}{}
von

\mathl{K[X_1 , \ldots ,  X_n]}{} nach $K$ in natürlicher Weise in Bijektion mit den Punkten aus dem
\definitionsverweis {affinen Raum}{}{}


\mavergleichskette

{\vergleichskette

{ { {\mathbb A}_{  K }^{  n   } }
}

{ = }{ K^n
}

{  }{
}

{    }{
}

{   }{
}

}
{}{}{,}}

\faktzusatz {und zwar entspricht dem Punkt

\mathl{\left(  a_1 , \ldots ,  a_n  \right)}{} der
\definitionsverweis {Einsetzungshomomorphismus}{}{}


\mathl{X_i \mapsto a_i}{.} Mit anderen Worten,


\mavergleichskettedisp

{\vergleichskette

{ K\!-\!\operatorname{Spek}\, \left(   K[X_1 , \ldots , X_n]   \right)
}

{ =} { { {\mathbb A}_{  K }^{  n   } }
}

{ } {
}

{ } {
}

{ } {
}

}
{}{}{.}}

\faktzusatz {}


}

{


Ein $K$-Algebrahomomorphismus ist stets durch ein $K$-Algebra-Er\-zeugendensystem festgelegt. D.h. die Werte an den Variablen $X_i$ legen einen $K$-Algebrahomomorphismus von

\mathl{K[X_1 , \ldots , X_n]}{} nach $K$ fest. Ein solcher Einsetzungshomomorphismus ist durch

\mathl{X_i \mapsto a_i}{} definiert. Zugleich ist hier jede Vorgabe von Werten

\mathl{\left(  a_1 , \ldots , a_n  \right)}{} erlaubt.


}


\inputbeispiel{}

{


Das $K$-Spektrum zur $K$-Algebra $K$ besteht einfach aus einem Punkt, und zwar ist die Identität
\maabb {} { K } { K
} {}
der einzige $K$-Algebrahomo\-morphismus von $K$ nach $K$. Es gibt im Allgemeinen weitere Körperautomorphismen auf $K$, doch diese sind keine $K$-Algebrahomomorphismen.


}


Entscheidend ist nun der folgende Satz, der eine bijektive Beziehung zwischen dem $K$-Spektrum von $R$ und dem Nullstellengebilde stiftet, das von einer Restklassendarstellung von $R$ herrührt.


\inputfaktbeweis

{Endlich erzeugte K-Algebren/K-Spektrum/Isomorph zu Einbettung/Fakt}

{Satz}

{}

{


\faktsituation {}

\faktvoraussetzung {Es sei $K$ ein
\definitionsverweis {Körper}{}{}
und sei $R$ eine
\definitionsverweis {endlich erzeugte}{}{}
kommutative
$K$-\definitionsverweis {Algebra}{}{}
mit
$K$-\definitionsverweis {Spektrum}{}{}


\mathl{K\!-\!\operatorname{Spek}\, \left(   R   \right)}{.} Es sei


\mavergleichskette

{\vergleichskette

{ R
}

{ = }{ K[X_1 , \ldots ,  X_n]/{\mathfrak a}
}

{  }{
}

{    }{
}

{   }{
}

}
{}{}{}
eine
\definitionsverweis {Restklassendarstellung}{}{}
von $R$ mit dem zugehörigen Restklassenhomomorphismus
\maabbdisp {\varphi} { K[X_1 , \ldots ,  X_n] } { R
} {}
und dem
\definitionsverweis {Nullstellengebilde}{}{}


\mavergleichskette

{\vergleichskette

{ V( {\mathfrak a} )
}

{ \subseteq }{ { {\mathbb A}_{  K }^{  n   } }
}

{  }{
}

{    }{
}

{   }{
}

}
{}{}{.}}

\faktfolgerung {Dann stiftet die
\definitionsverweis {Abbildung}{}{}
\maabbeledisp {} { K\!-\!\operatorname{Spek}\, \left(   R   \right)  } { { {\mathbb A}_{  K }^{  n   } }
} { P } { P \circ \varphi
} {}
eine
\definitionsverweis {Bijektion}{}{}
zwischen

\mathl{K\!-\!\operatorname{Spek}\, \left(   R   \right)}{} und

\mathl{V( {\mathfrak a} )}{,} die bezüglich der
\definitionsverweis {Zariski-Topologie}{}{}
ein
\definitionsverweis {Homöomorphismus}{}{}
ist.}

\faktzusatz {}

\faktzusatz {}


}

{


Zunächst ist die angegebene Abbildung wohldefiniert, da die Hintereinanderschaltung


\mathdisp {P \circ \varphi : K[X_1 , \ldots , X_n] \stackrel{\varphi}{\longrightarrow} K[X_1 , \ldots , X_n]/ {\mathfrak a}  \cong R

 \mathdisplaybruch \stackrel{P}{\longrightarrow} K} {  }


einen $K$-Algebrahomomorphismus vom Polynomring nach $K$ definiert, der nach
Lemma 12.3
der Einsetzungshomomorphismus zu

\mathl{(a_1 , \ldots , a_n)}{} ist und mit dem entsprechenden Punkt des affinen Raumes identifiziert werden kann
\zusatzklammer {und zwar ist


\mavergleichskettek

{\vergleichskettek

{ a_i
}

{ = }{ P(\varphi(X_i))
}

{  }{
}

{    }{
}

{   }{
}

}
{}{}{}} {} {.}


Da der Homomorphismus

\mathl{P \circ \varphi}{} durch $R$ faktorisiert, wird das Ideal ${\mathfrak a}$ auf $0$ abgebildet. D.h. der Bildpunkt


\mavergleichskette

{\vergleichskette

{ P \circ \varphi
}

{ = }{ (a_1 , \ldots , a_n)
}

{  }{
}

{    }{
}

{   }{
}

}
{}{}{}
liegt in

\mathl{V( {\mathfrak a} )}{,} und es liegt eine Abbildung
\maabbeledisp {} { K\!-\!\operatorname{Spek}\, \left(   R   \right) } { V({\mathfrak a} ) \subseteq { {\mathbb A}_{  K }^{  n   } }
} { P } { P \circ \varphi
} {}
vor, die wir als bijektiv nachweisen müssen.


Es seien dazu


\mavergleichskette

{\vergleichskette

{ P_1, P_2
}

{ \in }{ K\!-\!\operatorname{Spek}\, \left(  R  \right)
}

{  }{
}

{    }{
}

{   }{
}

}
{}{}{}
zwei verschiedene Punkte. Es liegen also zwei verschiedene $K$-Algebrahomomorphismen vor, und da ein $K$-Algebraho\-momorphismus auf einem $K$-Algebra-Erzeugendensystem festgelegt ist, müssen sich die beiden auf mindestens einer Variablen unterscheiden. Dann ist aber auch der Wert der zugehörigen Koordinate verschieden, d.h.

\mathl{P_1 \circ \varphi  \neq P_2 \circ \varphi}{,} und die Abbildung ist injektiv.


Zur Surjektivität sei ein Punkt


\mavergleichskette

{\vergleichskette

{ (a_1 , \ldots ,  a_n)
}

{ \in }{ V( {\mathfrak a} )
}

{  }{
}

{    }{
}

{   }{
}

}
{}{}{}
vorgegeben. Der zugehörige $K$-Algebrahomomorphismus
\maabbeledisp {} { K[X_1 , \ldots ,  X_n] } { K
} { X_i } { a_i
} {,}
annulliert daher jedes


\mavergleichskette

{\vergleichskette

{ F
}

{ \in }{ {\mathfrak a}
}

{  }{
}

{    }{
}

{   }{
}

}
{}{}{,}
sodass dieser Ringhomomorphismus durch

\mathl{K[X_1 , \ldots , X_n]/ {\mathfrak a}}{} faktorisiert. Dieser Ringhomomorphismus ist das gesuchte Urbild aus

\mathl{K\!-\!\operatorname{Spek}\, \left(  R  \right)}{.}


Zur Topologie muss man einfach nur beachten, dass für


\mavergleichskette

{\vergleichskette

{ G
}

{ \in }{ R
}

{  }{
}

{    }{
}

{   }{
}

}
{}{}{}
und ein Urbild


\mavergleichskette

{\vergleichskette

{ \tilde{G}
}

{ \in }{ K[X_1 , \ldots , X_n]
}

{  }{
}

{    }{
}

{   }{
}

}
{}{}{}
und einen Punkt


\mavergleichskette

{\vergleichskette

{ P
}

{ \in }{ K\!-\!\operatorname{Spek}\, \left(  R  \right)
}

{  }{
}

{    }{
}

{   }{
}

}
{}{}{}
mit Bildpunkt


\mavergleichskette

{\vergleichskette

{ \tilde{P}
}

{ = }{ P \circ \varphi
}

{ \in }{ V( {\mathfrak a} )
}

{    }{
}

{   }{
}

}
{}{}{}
gilt:


\mavergleichskettedisp

{\vergleichskette

{ G(P)
}

{ =} { P(G)
}

{ =} { P(\varphi(\tilde{G}))
}

{ =} { ( P \circ \varphi ) (\tilde{G})
}

{ =} { \tilde{G} (\tilde{P})
}

}
{}{}{,}
sodass auch die Nullstellen übereinstimmen.


}


Dieser Satz besagt also, dass man jedes $K$-Spektrum einer endlich erzeugten $K$-Algebra $R$ mit einer Zariski-abgeschlossenen Menge eines

\mathl{{ {\mathbb A}_{ K }^{ n  } }}{} identifizieren kann. Man spricht von einer \stichwort {abgeschlossenen Einbettung} {.}


\inputfaktbeweis

{Endlich erzeugte K-Algebren/Nullenstellengebilde zu verschiedenen Restklassendarstellungen sind isomorph/über K-Spektrum/Fakt}

{Korollar}

{}

{


\faktsituation {}

\faktvoraussetzung {Es sei $K$ ein
\definitionsverweis {Körper}{}{}
und $R$ eine
\definitionsverweis {endlich erzeugte}{}{}
kommutative
$K$-\definitionsverweis {Algebra}{}{}
mit zwei
\definitionsverweis {Restklassendarstellungen}{}{}


\mathdisp {R \cong K[X_1 , \ldots , X_n]/ {\mathfrak a} \text{   und } R \cong K[X_1 , \ldots , X_m]/{\mathfrak b}} {  }


mit zugehörigen
\definitionsverweis {Nullstellengebilden}{}{}


\mavergleichskette

{\vergleichskette

{ V({\mathfrak a})
}

{ \subseteq }{ { {\mathbb A}_{  K }^{  n   } }
}

{  }{
}

{    }{
}

{   }{
}

}
{}{}{}
und


\mavergleichskette

{\vergleichskette

{ V({\mathfrak b})
}

{ \subseteq }{ { {\mathbb A}_{  K }^{  m   } }
}

{  }{
}

{    }{
}

{   }{
}

}
{}{}{.}}

\faktfolgerung {Dann sind die beiden Nullstellengebilde
\mathkor {} {V({\mathfrak a})} {und} {V({\mathfrak b})} {}
mit ihrer induzierten
\definitionsverweis {Zariski-Topologie}{}{}
\definitionsverweis {homöomorph}{}{}
zueinander.}

\faktzusatz {}

\faktzusatz {}


}

{


Nach
Satz 12.5
sind beide Nullstellengebilde homöomorph zu

\mathl{K\!-\!\operatorname{Spek}\, \left(   R   \right)}{,} sodass sie auch untereinander homöomorph sein müssen.


}


Wenn $R$ der Nullring ist, so ist das $K$-Spektrum davon leer. Wenn $K$ nicht algebraisch abgeschlossen ist, so kann das Spektrum auch zu anderen Ringen leer sein. Wenn aber $K$ ein algebraisch abgeschlossener Körper ist, so ist bei


\mavergleichskette

{\vergleichskette

{ R
}

{ \neq }{ 0
}

{  }{
}

{    }{
}

{   }{
}

}
{}{}{}
das Spektrum auch nicht leer. In der Tat gilt unter dieser Voraussetzung wieder ein Hilbertscher Nullstellensatz, siehe
Aufgabe 12.10.


  \zwischenueberschrift{Das K-Spektrum als Funktor}


\inputfaktbeweis

{Affine Varietäten/K-Spektren als Funktor/Fakt}

{Satz}

{}

{


\faktsituation {}

\faktvoraussetzung {Es sei $K$ ein
\definitionsverweis {Körper}{}{}
und seien $R$ und $S$ kommutative
$K$-\definitionsverweis {Algebren von endlichem Typ}{}{.}
Es sei
\maabb {\varphi} { R } { S
} {}
ein
$K$-\definitionsverweis {Algebrahomomorphismus}{}{.}}

\faktfolgerung {Dann induziert dies eine Abbildung
\maabbeledisp {\varphi^*} { K\!-\!\operatorname{Spek}\, \left(   S   \right) } { K\!-\!\operatorname{Spek}\, \left(   R   \right)
} { P } { P \circ \varphi
} {.}
Diese Abbildung ist
\definitionsverweis {stetig}{}{}
bezüglich der
\definitionsverweis {Zariski-Topologie}{}{.}}

\faktzusatz {}

\faktzusatz {}


}

{


Die Existenz der Abbildung ist klar, dem $K$-Algebrahomomorphis\-mus
\maabbdisp {P} { S } { K
} {}
wird einfach die Hintereinanderschaltung


\mathdisp {R \stackrel{\varphi}{\longrightarrow} S \stackrel{P}{\longrightarrow} K} {  }


zugeordnet. Das Urbild der offenen Menge


\mavergleichskette

{\vergleichskette

{ D(f)
}

{ \subseteq }{ K\!-\!\operatorname{Spek}\, \left(  R  \right)
}

{  }{
}

{    }{
}

{   }{
}

}
{}{}{}
ist dabei


\mavergleichskettealignhandlinks

{\vergleichskettealignhandlinks

{ (\varphi^*)^{-1} (D(f) )
}

{ =} { \left\{ P \in K\!-\!\operatorname{Spek}\, \left(  S  \right)   \mid  \varphi^*(P) \in D(f)         \right\}
}

{ =} { \left\{ P \in K\!-\!\operatorname{Spek}\, \left(  S  \right)   \mid   P \circ \varphi  \in D(f)         \right\}
}

{ =} { \left\{ P \in K\!-\!\operatorname{Spek}\, \left(  S  \right)   \mid   ( P \circ \varphi) (f) \neq 0         \right\}
}

{ =} { \left\{ P \in K\!-\!\operatorname{Spek}\, \left(  S  \right)   \mid   P ( \varphi(f)) \neq 0         \right\}
}

}
{

\vergleichskettefortsetzungalign

{ =} { D( \varphi(f))
}

{ } {}

{ } {}

{ } {}

}
{}{.}
Daher sind generell Urbilder von offenen Mengen wieder offen und die Abbildung ist stetig.


}


Die in
Satz 12.7
eingeführte Abbildung $\varphi^*$ nennt man die \stichwort {Spektrumsabbildung} {} zu $\varphi$.


\inputfaktbeweis

{Endlich erzeugte K-Algebren/K-Spektren als Funktor/Verschiedene Homomorphismen/Fakt}

{Proposition}

{}

{


Es sei $K$ ein
\definitionsverweis {Körper}{}{}
und zu einem
$K$-\definitionsverweis {Algebrahomo\-morphismus}{}{}
\maabb {\varphi} { R } { S
} {}
zwischen
$K$-\definitionsverweis {Algebren von endlichem Typ}{}{}
sei $\varphi^*$ die zugehörige
\definitionsverweis {Spektrumsabbildung}{}{.}
Dann gelten folgende Aussagen.
\aufzaehlungfuenf{Zu einem $K$-Algebrahomomorphismus
\maabb {P} { R } { K
} {}
ist die induzierte Spektrumsabbildung $P^*$ einfach die Abbildung, die dem einzigen Punkt


\mavergleichskette

{\vergleichskette

{ \{ \operatorname{id}\}
}

{ = }{ K\!-\!\operatorname{Spek}\, \left(   K   \right)
}

{  }{
}

{    }{
}

{   }{
}

}
{}{}{}
den Punkt


\mavergleichskette

{\vergleichskette

{ P
}

{ \in }{ K\!-\!\operatorname{Spek}\, \left(   R   \right)
}

{  }{
}

{    }{
}

{   }{
}

}
{}{}{}
zuordnet.
}{Der durch ein Element


\mavergleichskette

{\vergleichskette

{ F
}

{ \in }{ R
}

{  }{
}

{    }{
}

{   }{
}

}
{}{}{}
definierte
\definitionsverweis {Einsetzungshomomorphismus}{}{}
\maabbeledisp {\varphi} { K[T] } { R
} { T } { F
} {,}
induziert die Spektrumsabbildung
\maabbeledisp {\varphi^*} { K\!-\!\operatorname{Spek}\, \left(   R   \right)  } { K\!-\!\operatorname{Spek}\, \left(  K[T]   \right) = {\mathbb A}^{1}_{ K }
} { P } { F(P)
} {.}
}{Zu einer
\definitionsverweis {surjektiven}{}{}
Abbildung
\maabb {\varphi} { R } { S
} {}
von $K$-Algebren von endlichem Typ ist die zugehörige Spektrumsabbildung
\maabbdisp {\varphi^*} { K\!-\!\operatorname{Spek}\, \left(   S   \right) } { K\!-\!\operatorname{Spek}\, \left(   R   \right)
} {}
eine
\definitionsverweis {abgeschlossene Einbettung}{}{,}
und zwar ist das Bild gleich

\mathl{V(\ker (\varphi))}{.}
}{Die zu einer surjektiven Abbildung
\maabb {} { K[X_1 , \ldots , X_n] } { S
} {}
gehörende Spektrumsabbildung
\maabbdisp {\varphi^*} { K\!-\!\operatorname{Spek}\, \left(   S   \right) } { K\!-\!\operatorname{Spek}\, \left(   K[X_1 , \ldots , X_n]   \right) \cong { {\mathbb A}_{  K }^{  n   } }
} {}
stimmt mit der in
Satz 12.5
definierten Abbildung überein.
}{Es seien


\mavergleichskette

{\vergleichskette

{ F_i
}

{ \in }{ K[X_1 , \ldots , X_n]
}

{  }{
}

{    }{
}

{   }{
}

}
{}{}{}
für


\mavergleichskette

{\vergleichskette

{ i
}

{ = }{ 1 , \ldots , m
}

{  }{
}

{    }{
}

{   }{
}

}
{}{}{}
und es sei
\maabbeledisp {\varphi} { K[Y_1 , \ldots , Y_m] } { K[X_1 , \ldots , X_n]
} { Y_i } { F_i
} {,}
der zugehörige
\definitionsverweis {Einsetzungshomomorphismus}{}{.}
Dann stimmt die
\definitionsverweis {Spektrumsabbildung}{}{}
\maabbdisp {\varphi^*} { { {\mathbb A}_{  K }^{  n   } } = K\!-\!\operatorname{Spek}\, \left(   K[X_1 , \ldots , X_n]   \right)  } { { {\mathbb A}_{  K }^{  m   } } = K\!-\!\operatorname{Spek}\, \left(   K[Y_1 , \ldots , Y_m]   \right)
} {}
\zusatzklammer {über die Identifizierung aus
Lemma 12.3} {} {}
mit der direkten
\definitionsverweis {polynomialen Abbildung}{}{}


\mathdisp {\left(  x_1 , \ldots , x_n  \right) \longmapsto \left(  F_1 \left(  x_1 , \ldots , x_n  \right) , \ldots , F_m \left(  x_1 , \ldots , x_n  \right)  \right)} {  }


überein.
}


}

{


(1) Dies folgt aus


\mavergleichskette

{\vergleichskette

{
\operatorname{Id} \circ P
}

{ = }{ P
}

{  }{
}

{    }{
}

{   }{
}

}
{}{}{.}


(2) Unter der hintereinandergeschalteten Abbildung


\mathdisp {K[T] \stackrel{\varphi }{\longrightarrow} R \stackrel{P}{\longrightarrow} K} {  }


wird $T$ auf


\mavergleichskette

{\vergleichskette

{ P(F)
}

{ = }{ F(P)
}

{  }{
}

{    }{
}

{   }{
}

}
{}{}{}
geschickt.


(3) beruht auf ähnlichen Betrachtungen, wie sie im Beweis zu
Satz 12.5
durchgeführt wurden. Das zeigt auch (4). Zu (5) siehe
Aufgabe 12.18.


}


Die unter (2) formulierte Aussage besagt insbesondere, dass man die Elemente des Ringes $R$ als Funktionen auf dem $K$-Spektrum

\mathl{K\!-\!\operatorname{Spek}\, \left(   R   \right)}{} nach ${\mathbb A}^{1}_{ K }$ auffassen kann. Wir haben also ein geometrisches Objekt eingeführt, mit dem man Ringelemente als Funktionen realisieren kann.


  \zwischenueberschrift{Weitere Eigenschaften des K-Spektrums}


\inputfaktbeweis

{Affine Varietäten/K-Spektren/Polynomring und affine Gerade/Fakt}

{Lemma}

{}

{


\faktsituation {}

\faktvoraussetzung {Es sei $K$ ein Körper und $R$ eine endlich erzeugte kommutative $K$-Algebra.}

\faktfolgerung {Dann ist


\mavergleichskettedisp

{\vergleichskette

{ K\!-\!\operatorname{Spek}\, \left(  R[X]   \right)
}

{ \cong} { K\!-\!\operatorname{Spek}\, \left(  R   \right)  \times {\mathbb A}^{1}_{ K }
}

{ } {
}

{ } {
}

{ } {
}

}
{}{}{.}}

\faktzusatz {}

\faktzusatz {}


}

{


Ein $K$-Algebrahomomorphismus
\maabb {} {R[T] } { K
} {}
induziert einen $K$-Al\-gebrahomomorphismus
\maabb {} { R } { K
} {,}
und zugleich wird $T$ auf ein bestimmtes Element


\mavergleichskette

{\vergleichskette

{ a
}

{ \in }{ K
}

{  }{
}

{    }{
}

{   }{
}

}
{}{}{}
abgebildet. Diese Daten definieren aber auch einen eindeutig bestimmten $K$-Algebrahomomorphismus
\maabb {} {R[T] } { K
} {.}


}


Achtung: Die vorstehende Aussage liefert nur eine natürliche Bijektion auf der Punktebene. Würde man die Produktmenge rechts mit der Produkttopologie versehen, so würde hier keine Homöomorphie mit der Zariski-Topologie links vorliegen. Insbesondere ist


\mavergleichskette

{\vergleichskette

{ {\mathbb A}^{2}_{K}
}

{ = }{ {\mathbb A}^{1}_{K} \times {\mathbb A}^{1}_{K}
}

{  }{
}

{    }{
}

{   }{
}

}
{}{}{,}
aber die Zariski-Topologie der affinen Ebene ist nicht die Produkt-Topologie der affinen Geraden mit sich selbst.


\inputbemerkung

{}

{
Sind


\mavergleichskette

{\vergleichskette

{ X
}

{ = }{ K\!-\!\operatorname{Spek}\, \left(   R   \right)
}

{  }{
}

{    }{
}

{   }{
}

}
{}{}{}
und


\mavergleichskette

{\vergleichskette

{ Y
}

{ = }{ K\!-\!\operatorname{Spek}\, \left(   S   \right)
}

{  }{
}

{    }{
}

{   }{
}

}
{}{}{,}
so  lässt sich die Produktmenge

\mathl{X \times Y}{} ebenfalls als $K$-Spektrum einer $K$-Algebra darstellen, und zwar ist


\mavergleichskettedisp

{\vergleichskette

{ X \times Y
}

{ \cong} { K\!-\!\operatorname{Spek}\, \left(   R \otimes_{ K } S  \right)
}

{ } {
}

{ } {
}

{ } {
}

}
{}{}{,}
wobei $\otimes$ das Tensorprodukt bezeichnet. Wir werden darauf nicht im Einzelnen eingehen. Um aber doch ein Gefühl dafür zu geben, betrachten wir


\mavergleichskette

{\vergleichskette

{ R
}

{ = }{ K[X_1 , \ldots , X_n]/ {\mathfrak a}
}

{  }{
}

{    }{
}

{   }{
}

}
{}{}{}
und


\mavergleichskette

{\vergleichskette

{ S
}

{ = }{ K[Y_1 , \ldots ,  Y_m]/  {\mathfrak b}
}

{  }{
}

{    }{
}

{   }{
}

}
{}{}{.}
Dann ist


\mavergleichskettedisp

{\vergleichskette

{ R \otimes_{ K } S
}

{ \cong} { K[X_1  , \ldots ,  X_n,Y_1 , \ldots , Y_m]/({\mathfrak a} +{\mathfrak b} )
}

{ } {
}

{ } {
}

{ } {
}

}
{}{}{}
\zusatzklammer {bei dieser ad hoc Definition ist nicht klar, dass sie unabhängig von den Darstellungen als Restklassenring ist} {} {.}
}
